\documentclass[10pt,journal,compsoc]{IEEEtran}
\usepackage[utf8]{inputenc}
\usepackage[T1]{fontenc}
\usepackage{url}
\usepackage[hidelinks]{hyperref}
\usepackage{cite}
\usepackage{amsmath,amssymb,amsfonts}
\usepackage{algorithmic}
\usepackage{graphicx}
\usepackage{textcomp}
\usepackage{xcolor}
\usepackage{tabularx}
\usepackage{listings}
\usepackage{array}

\lstset{
  basicstyle=\small\ttfamily,
  breaklines=true,
  frame=single,
  columns=fullflexible
}

\begin{document}

\title{Module 004 Stdp Causal}
\author{Bryan Alexander Buchorn \\ Independent Researcher \\ bryan.alexander@buchorn.com}
\markboth{CEREBRUM v1.2.4 Hardened Edition · March 2026}{}

\maketitle

\begin{abstract}
This module forms a critical component of the CEREBRUM framework, a community-structured graph attention engine for Knowledge Graph reasoning. It focuses on Advanced Graph Attention Analysis.
\end{abstract}

\begin{IEEEkeywords}
Knowledge Graph, Graph Attention, DSCF, CSA, Hierarchical Reasoning.
\end{IEEEkeywords}

\section{Autonomous Causal Discovery via Spike-Timing-Dependent Plasticity in Knowledge Streams}

\textbf{Authors}: Bryan Alexander Buchorn  
\textbf{Affili\-ations}: Independent Researcher, Las Vegas, NV, USA  
\textbf{Status}: v1.2.0 (Hardened Enterprise)  
\textbf{Date}: March 2026

---

\subsubsection{Abstract}
Inferring causal relat\-ionships from unstr\-uctured, high-velocity event streams is a major challenge in unsup\-ervised learning. We propose a novel method for autonomous causal discovery by adapting the biological mechanism of \textbf{Spike-Timing-Dependent Plasticity (STDP)} to temporal Knowledge Graph triples. By treating entity mentions as "spikes" and analyzing their relative timing across a sliding window, our engine infers directional \texttt{CAUSES} relat\-ionships. We define a mathe\-matically rigorous weighting rule based on the Bi \& Poo \cite{bipoo1998} model and introduce \textbf{Lazy Decay}, an $O(1)$ optim\-ization that allows the engine to scale to enterprise-level event throughput by applying geometric decay only upon record access. Benchmark results demonstrate that the v1.2.0 engine maintains constant sub-millisecond latency per event regardless of the number of accumulated causal pairs, repre\-senting a critical break\-through for real-time industrial causal monitoring.

\subsubsection{1. Introd\-uction}
Traditional causal inference relies on static datasets and intensive count\-erfactual analysis. In streaming envir\-onments—such as IoT telemetry, cyber\-security logs, or financial tickers—causality must be discovered "on the fly." We posit that the temporal order and proximity of events provide a sufficient signal for preliminary causal discr\-etization when governed by biological plasticity rules.

\subsubsection{2. Methodology}

\paragraph{2.1 The STDP Weighting Rule}
For any pair of entity spikes $(u, v)$ where $u$ occurs at $t_{pre}$ and $v$ at $t_{post}$, the causal weight $w_{uv}$ is updated based on the interval $\Delta t = t_{post} - t_{pre}$:
\begin{itemize}
\item \textbf{LTP (Potent\-iation)}: $\Delta w_{uv} = A_+ \cdot e^{-\Delta t / \tau_+}$ if $\Delta t > 0$.
\item \textbf{LTD (Depression)}: $\Delta w_{uv} = -A_- \cdot e^{-|\Delta t| / \tau_-}$ if $\Delta t < 0$ or if $v$ spikes without a preceding $u$.

\end{itemize}
\paragraph{2.2 Signif\-icance Filtering}
To distinguish causal signal from rhythmic or stochastic noise, we apply a four-stage filter:
\begin{enumerate}
\item \textbf{Weight Threshold}: $w_{uv} \geq w_{threshold}$.
\item \textbf{Pairing Count}: $n \geq n_{min}$ (minimum evidence).
\item \textbf{Temporal Span}: Minimum wall-clock duration between first and last pairing (Phase 19).
\item \textbf{Uniformity}: A $\chi^2$ test on inter-spike intervals to reject burst-driven artifacts (Phase 19).

\end{enumerate}
\subsubsection{3. Scalability: Lazy Decay}
In standard imple\-mentations, decaying $N$ weights is $O(N)$. We introduce \textbf{Lazy Decay}, which computes the accumulated decay for a specific pair only upon access:
\begin{equation}
w'_{uv} = w_{uv} \cdot \lambda^{(T - t_{last})}
\end{equation}
where $T$ is the global step and $t_{last}$ is the pair's last update. This reduces the complexity of causal maintenance from $O(N)$ to $O(1)$ per event.

\subsubsection{4. Conclusion}
The STDP Causal Engine provides a scalable, unsup\-ervised framework for real-time causal discovery. By grounding graph evolution in the temporal dynamics of the data stream, it enables autonomous reasoning engines to identify and follow causal chains as they emerge.

---
\textbf{References}
\begin{enumerate}
\item Bi, G. Q., \& Poo, M. M. (1998). Synaptic modif\-ications in cultured hippocampal neurons. Journal of Neuros\-cience.
\item Markram, H., et al. (1997). Regulation of synaptic efficacy by predictions of spike timing. Science.
\item Pearl, J. (2009). Causality: Models, Reasoning, and Inference. Cambridge University Press.
\item Spirtes, P., Glymour, C. N., \& Scheines, R. (2000). Causation, Prediction, and Search. MIT Press.
\item Sjöström, J., \& Gerstner, W. (2010). Spike-timing dependent plasticity. Schola\-rpedia.
\item Song, S., Miller, K. D., \& Abbott, L. F. (2000). Competitive Hebbian learning through spike-timing-dependent synaptic plasticity. Nature Neuros\-cience.
\item Buchorn, B. A., \& Sonnet, C. (2026). Lazy STDP Weight Decay in CEREBRUM. SPEC\_004.md.

\end{enumerate}

\bibliographystyle{IEEEtran}
\bibliography{../../docs/latex/references}

\end{document}
